\subsection{Lifting the columns of a path}
\label{sec:column-lifts}
\hypertarget{source-column-lifts}{}

To obtain \eqref{eq:raw-tail-required}, one needs integral polynomials realizing the retained columns with control of the terms beyond their initial symbols. Fix a level of
a block as in Section~\ref{sec:segment-window}. The chain at this
level consists of the reduced vector and covector degrees in
dimension \(r=k-1\). Its prefixes have the first-row bounds
\(\lfloor K_l-\mu\rfloor\), where
\(\mu=M_0+nA_1\). We allow \(M_0=0\) or \(12c\).
Powers of the pairings \(E_i,E_+\), scalar decorations, and the
remaining endpoint covector monomials multiply this chain space.

\begin{paramdefinition}[Standard words of a Pieri path]\label{hyp:standard-path-words}
For the chain alone take a Pieri path as in
Section~\ref{sec:retained-rank}, with successive dominant weights
\(\lambda(t)\) starting at zero. For each integer column \(b\)
put
\[
 h_b(t)=\#\{j:\lambda_j(t)\ge b\}.
\]

In a positive column keep an exterior tensor of degree \(h_b\)
in the vector representation. For \(b\le0\) use the dual
representation and degree \(r-h_b\). Begin with 1. When the degree
rises, wedge on the right with the letter, and when it falls,
contract on the right with the letter. Monotonicity and interlacing
give only changes by one. In the dual exterior tensor a covector is
wedged and a vector contracted. When the height stays fixed, do
nothing. In the ordered subset basis, inserting or deleting index
\(j\) contributes the sign
\[
 (-1)^{\#\{u\in J:u>j\}},
\]
with \(J\)
the subset after omitting the inserted or deleted index.

Choose an output subset in each column. Complement these subsets
in the nonpositive columns, then require the resulting columns,
sorted increasingly, to form a weakly increasing tableau along
each row. Only finitely many columns act, since far to the left
the complemented column is full and far to the right it is empty.
The product of the selected exterior coordinates is the standard
word of the path and tableau. It is integral and of exactly the
specified multidegree, since a strip uses as many distinct columns
as its size.
\end{paramdefinition}

\begin{lemma}[Integral Pieri path words]
\label{lem:integral-path-basis}
The words of Definition~\ref{hyp:standard-path-words}
form an integral basis in the specified multidegree. If \(\mathcal W_{\rm pre}\) is the subset of words satisfying the prefix conditions and \(U_{\rm pre}\) the corresponding rational subspace of the lattice \(L\) tensored with \(\mathbb Q\), then
\[
 L\cap U_{\rm pre}=\bigoplus_{w\in\mathcal W_{\rm pre}}\mathbb Z w.
\]

\end{lemma}

\begin{proof}
\proofstep{The triangularity of the path basis.}
Here is a basis proof which also explains the importance of using
integral words. In a putative linear relation over a field insert
after each letter \(t\) an independent formal diagonal matrix
\[
 D_t=\operatorname{diag}(T_{t1},\ldots,T_{tr})
\]
acting on every
live exterior tensor, and after all letters insert one generic
invertible matrix \(g\). The relation still holds over the field
of rational functions. Indeed wedge and contraction commute with
a simultaneous change of basis, using its contragredient on dual
tensors. Moving the matrices past the subsequent operations to
the inputs changes each letter by the product of matrices following
it. This is one substitution, common to all paths.

To find the largest exponent vector of \(D_t\) on a fixed
path, take the initial \(h_b\) indices in a positive column \(b\)
and the last \(r-h_b\) indices in a dual column. If \(I_b(t)\)
is the chosen set and \(\varepsilon_b=1,-1\) for the two column
signs, the exponent of \(T_{ti}\) in the product is
\[
 \lambda_i(t)=\sum_b \varepsilon_b\,\mathbf1_{i\in I_b(t)} .
\]
Each chosen exponent vector is lexicographically maximal in its
column, and a strict decrease in any column strictly decreases
the sum in that order. Thus the maximal term forces every
choice \(I_b(t)\). These choices are compatible at consecutive
steps.

In the one changing column the old and new sets differ
by adding or deleting a single index \(i\); the corresponding
matrix entry of wedge or contraction by the input letter is
its \(i\)-th coordinate times a sign. Hence the coefficient
at the tuple of all these maxima is a signed monomial in the
inputs times the terminal matrix coefficient in \(g\). Index the insertions and deletions in the selected columns by \(a\), and let \(v_{a,i_a}\) be the input coordinate at occurrence \(a\). Their path factor is
\[
 \varepsilon\prod_a v_{a,i_a},\qquad \varepsilon\in\{1,-1\}.
\]
All dependence on the terminal tableau lies in the matrix coefficient. The matrix coefficient is a
product of minors with initial rows.

For a dual column, the
complementary minor identity introduces a sign and one inverse
determinant. After a common determinant power clears these for a
fixed path, the result consists of ordinary flag minors with
the tableau columns. In lexicographic order on the matrix entries,
row by row with smaller column index preferred, a flag minor with sorted column indices \(i_1<\cdots<i_h\) has leading monomial
\[
 \operatorname{in}\det(g_{a,i_b})_{a,b=1}^h=\prod_{a=1}^h g_{a,i_a}.
\]
The determinant signs from the dual columns change only the overall sign. The diagonal
monomial recovers each sorted tableau row. Consequently these
matrix coefficients are independent for each path, in every
characteristic. Choose the lexicographically largest sequence of
weights in the relation, comparing the diagonal matrices in
temporal order. The coefficient just computed forces all the
coefficients for that sequence to vanish. This proves independence.

\proofstep{The rank and prefix operators.}
Let \(p_\lambda\) count paths ending at \(\lambda\), and let
\(d_\lambda\) count its semistandard tableaux. For integer weights
the tableau can be shifted by a whole determinant power, which
does not change \(d_\lambda\). Pieri gives
\[
 \sum_\lambda p_\lambda d_\lambda
       =\prod_j\binom{r+o_j-1}{r-1},
\]
the monomial rank for the prescribed letter degrees \(o_j\).
The integer coefficient matrix is therefore square. Write \(A\) for this matrix. Independence in each residue characteristic gives
\[
 \det A\not\equiv0\pmod p\quad(\text{every prime }p),\qquad \det A\in\mathbb Z^\times=\{1,-1\}.
\]


At a prefix \(t\), varying the first \(t\) letters by a common
matrix factors their contribution through the tensor of the exterior
powers at that cut. In each torus coordinate its largest possible
weight is at most \(\lambda_1(t)\), by the same signed subset
calculation. The word therefore belongs to the first-row filtered
space at that bound. Over characteristic zero the multiplicity-one Pieri decompositions give
\[
 \dim\bigcap_t\{\text{prefix cutoff at }t\}=\#\{\text{basis words whose paths satisfy every cutoff}\}.
\]
Indeed one
may first split a prefix tensor into irreducibles and then tensor
each summand with the next letter. The later isotypic projections
preserve these summands. The independent allowed words consequently
span that intersection. Since they form a subset of an integral
basis, they also span its intersection with the coefficient lattice.
Expanding at each level and using the monomial basis for the
endpoint covector gives the desired generators of the integral
symbol space.

\end{proof}

\paragraph{Putting back the first row.}
Write \(q_j=p^{G_j}\) for the radix at the start of segment \(j\),
\(B=\lfloor(a-E_0)/q_j\rfloor\), and \(z_i=A_i-w_i\) or
\(A_i-o_i\) according as \(i<\sigma\) or \(i\ge\sigma\).
Prepend an even vector letter. Enlarge the path by a new first
row, initially zero, which at that first letter rises to
\[
             b_0=\lfloor B-\mu\rfloor-\sum_{i<\sigma} h_i .
\]
Increase this row by \(z_i\) at the odd vector \(i\). Keep it
fixed at the covectors and at the final vector. The reduced
path is unchanged in the other rows, delayed by the prepended
step.

For every vector the first reduced row after the addition
must be at most the enlarged first row \emph{before} the addition.
At a tied vector and the final vector this first row before
addition is exactly \(\lfloor K_i-\mu\rfloor\) and
\(\lfloor K_s-\mu\rfloor\), respectively. The initial \(h_i\)
are all accounted for in \(b_0\).

For the initial vectors the reduced first row is at most
\(\sum_{l\le i}w_l\). Put \(\epsilon_0=N_1\mathcal L/k\) and
\[
 e_0=(c-n)A_1.
\]
In the local size range
\(b'\ge k^3NQ^4cY,\; kb'\le a,\;N\ge2,\;Y\ge1\).
The floor estimate
\[
 B\ge a/(2Q)
\]
gives
\[
 \epsilon_0B\ge k^3Q^3 c,
\]
using \(\epsilon_0\ge1/k\).
Since \(2N_1+3\le k\) and \(k\ge5\),
\[
 \mu+N_1e_0\le12c+13kN_1^2Qc\le4\epsilon_0 B,\qquad
 \sum_{i<\sigma}h_i\le 2\epsilon_0B+N_1e_0 .
\]
In particular \(b_0\ge0\).

By \eqref{eq:coarse-prefix},
\(\sum_{l\le i} w_l\le(1+4/p^2)A_i\).
Using \((1+\Delta)A_i<B\), \(4/p^2\le\epsilon_0\),
\(40\epsilon_0\le\Delta<1/6\), we get
\[
 (1+4/p^2)A_i+6\epsilon_0 B\le B.
\]
Thus the required interlacing holds at every initial vector too.
The first row stays fixed at a removal, for which old dominance
and reduced interlacing suffice. This constructs a valid path in
dimension \(1+r\), including for an empty odd list.

\paragraph{Frobenius powers in exterior coordinates.}
For a column of the enlarged path perform the same exterior
recursion on the true letters, with the following insertions.
For successive used Frobenius indices \(s,s'\), the intervening operation on an exterior coordinate \(c_I\) is
\[
 c_I\longmapsto c_I^{\,p^{s'-s}}.
\]
Use the terminal index of the segment for the last such operation. Use increasing integer lifts of
the positions through wrap-around. The vector and covector at
one odd position have the same index. The result is integral.
At a used position \(s\) its degree in that letter is
\(p^{s_{\rm end}-s}\), and zero at unused positions. Positions
of the same orientation in a segment are distinct, including
the start and end by the two-segment hypothesis. Each true
letter degree in this column is therefore at most \(Q\).

Use coordinate 0 for the added line and the remaining coordinates
for the reduced tensor. We claim that the new column has order
at least
\begin{equation}\label{eq:lifted-path-column-order}
 d_b=\sum_{s:\,\text{reduced column height rises at }s}
                         p^{s_{\rm end}-s}.
\end{equation}
at every residue line. Its symbol after division by \(w^{d_b}\)
at the identity is the reduced column with every letter raised
coordinatewise to its radix, times the pure first-coordinate
power at the step at which the new first row crosses \(b\),
if it crosses at all. For a positive column include coordinate
0 in the output subset exactly when the first row reaches \(b\).
For a nonpositive column keep the dual output subset in the
remaining coordinates. Use the ordering already specified, so
right wedge and right contraction on ordinary coordinates have
the reduced signs.

\begin{lemma}[Order of a lifted path column]
\label{lem:column-graded-lift}
For the enlarged path at this block level, let \(C_\ell\) be its column at a residue line \(\ell\), with integral coefficient lattice \(L_\ell\). The order \(d_b\) in \eqref{eq:lifted-path-column-order} satisfies
\[
 C_\ell\in w^{d_b}L_\ell\qquad(\text{each residue line }\ell).
\]

At the identity line its symbol modulo \(w\) after this division is
the reduced column evaluated on the coordinatewise radix powers,
multiplied by the crossing scalar when the first row is crossed. The
symbol is homogeneous of normal degree \(d_b\).
\end{lemma}

\begin{proof}
\proofstep{The orders after the Frobenius powers are lifted.}
We give the order proof since taking Frobenius powers in a
moving frame requires two orders of precision. Call a subset
allowed in a positive column of positive height when it contains
0 and has that height, and at height zero when it is empty.
In a dual column the allowed subsets omit 0 and have the
complementary height. Measure the entries in the adapted frame
of the tested line after each step.

The induction asserts that
all are divisible by \(w^d\) and the entries outside the allowed
subsets by \(w^{d+1}\). Initially this holds for the scalar 1.
Between indices let \(P=p^{s'-s}\) and let \(L_s\) convert actual
exterior coordinates into adapted coordinates, using the
appropriate twisted exterior matrix. Write the adapted row as
\(w^d U\). After entrywise powering and conversion at \(s'\) it is
\[
         w^{Pd}\big(U^{[P]}+pV\big)
\]
for an integral polynomial row \(V\). Indeed expand
\((U L_s^{-1})^{[P]} L_{s'}\) modulo \(p\), where powering is
additive and
\[
 (L_s^{-1})^{[P]}\equiv L_{s'}^{-1}.
\]

The last congruence holds in the unramified Witt ring before
extension to \(O\), for inverse minors as well as minors.
At \(s'=s\) the identity is exact. If the index strictly
increases, the extra \(w\) in a disallowed entry becomes at least
\(w^2\) after the common factor is removed, since
\[
 p=w^2
\]

and \(P\ge2\).

In a positive column a first-row crossing wedges the free
vector coordinate 0 into an empty exterior tensor. For a reduced height increase the tensor already contains \(e_0\). Since
\[
 e_0\wedge(e_0\wedge v)=0,
\]
a surviving insertion must use an ordinary coordinate carrying \(w\).
A height decrease contracts an ordinary covector at no cost.
Contracting coordinate zero costs an extra \(w\) and is disallowed.
For a nonpositive column the complemented tableau always contains
the new first row. A decrease wedges an ordinary covector. An
increase contracts an ordinary vector at cost \(w\). These exhaust
the possible degree changes.

At either costly step a free coordinate acting on a disallowed entry could consume the spare power of \(w\). Every costly step is a vector
at a strictly later index than all preceding used steps, as the
vector precedes the covector at an equal index. After the common factor is removed, the disallowed term has order at least two. Paying the one power for a costly step therefore still leaves the required extra order:
\[
 \operatorname{ord}_w(\text{disallowed term})-1\ge2-1=1 .
\] If it is the first used step,
all the disallowed entries are zero. At a free step the extra
order is preserved directly. If \(d_{\rm old}\) is the common order before a step and
\(P\) is its intervening Frobenius radix, we have proved
\[
 d_{\rm new}=P d_{\rm old}+
       \mathbf1_{\{\text{reduced height rises at this step}\}} .
\]
Iterating from order zero weights each costly step by its
remaining radix \(p^{s_{\rm end}-s}\), exactly
\eqref{eq:lifted-path-column-order}. Inverting the final
integral frame change preserves this common divisibility.

\proofstep{Returning to the identity frame.}
At the identity the frame matrices are identities. In the
normalized leading part the steps just described act on the
ordinary exterior coordinates exactly as in the reduced word.
Crossing the first row occurs while the ordinary positive
column is empty, and contributes its scalar vector entry.
If \(P\) is the intervening Frobenius power, then in the residue field
\[
 \Big(\sum_I\varepsilon_I a_I b_I\Big)^P=\sum_I\varepsilon_I a_I^P b_I^P\qquad(\varepsilon_I\in\{1,-1\}).
\]
Thus applying Frobenius before the exterior operation gives the same residue. The resulting ordinary column is thus evaluated
on the stated radix powers.

The leading expression in
characteristic zero lifted this way is homogeneous of normal
degree \(d_b\), counting ordinary vector coordinates and scalar
covector coordinates with weight one. The argument proves a
congruence modulo \(w\) after normalization, as required.
Taking the product of the columns reproduces all the reduced
words and spends \(q_j b_0\) at the start and \(P_i z_i\)
in vector coordinate zero at each odd top.
\end{proof}
